% part: intuitionistic-logic
% chapter: soundness-completeness
% section: soundness-nd

\documentclass[../../../include/open-logic-section]{subfiles}

\begin{document}

\olfileid[hi]{int}{sc}{snd}

\olsection{प्राकृतिक निगमन की यथार्थता}

अब हम संबंधात्मक अर्थविज्ञान के सापेक्ष प्राकृतिक निगमन की यथार्थता सिद्ध
करेंगे; अर्थात दिखाएँगे कि यदि कोई !!{formula} मान्यताओं के किसी समुच्चय से
!!{derivable} है, तो उस मान्यता-समुच्चय से वह !!{formula} निष्पन्न होता है।

\begin{thm}[यथार्थता]
  \ollabel{thm:soundness}
  यदि $\Gamma \Proves !A$, तो $\Gamma \Entails !A$।
\end{thm}

\begin{proof}
  हम सिद्ध करते हैं कि यदि $\Gamma \Proves !A$, तो
  $\Gamma \Entails !A$। प्रमाण~$\Gamma$ से~$!A$ की !!{derivation} पर
  आगमन द्वारा है।
  
  \begin{enumerate}
  \item यदि !!{derivation} केवल मान्यता~$!A$ से बनी है, तो
    $!A \Proves !A$, और हमें $!A \Entails !A$ दिखाना है। मान लें
    $\mSat{M}{!A}[w]$। तब तुच्छतः $\mSat{M}{!A}[w]$।

  \item !!{derivation} $\Intro{\land}$ पर समाप्त होती है:
    \iftag{probAnd}{अभ्यास।}{!!{undischarged} मान्यताओं~$\Gamma$ से पूर्वाधार
      $!B$ की और !!{undischarged} मान्यताओं~$\Delta$ से~$!C$ की
      !!{derivation}याँ दिखाती हैं कि $\Gamma \Proves !B$ और
      $\Delta \Proves !C$। आगमन-परिकल्पना से $\Gamma \Entails !B$ और
      $\Delta \Entails !C$। हमें
      $\Gamma \cup \Delta \Entails !A \land !B$ दिखाना है, क्योंकि पूरी
      !!{derivation} की !!{undischarged} मान्यताएँ $\Gamma$ और $\Delta$ साथ
      मिलकर हैं। अतः मान लें $\mSat{M}{\Gamma \cup \Delta}[w]$। तब
      $\mSat{M}{\Gamma}[w]$ भी। चूँकि $\Gamma \Entails !B$,
      $\mSat{M}{!B}[w]$। इसी प्रकार $\mSat{M}{!C}[w]$। अतः
      $\mSat{M}{!B \land !C}[w]$।}

  \item !!{derivation} $\Elim{\land}$ पर समाप्त होती है:
    \iftag{probAnd}{अभ्यास।}{!!{undischarged} मान्यताओं $\Gamma$ से पूर्वाधार
        $!B \land !C$ की !!{derivation} दिखाती है कि
        $\Gamma \Proves !B \land !C$। आगमन-परिकल्पना से
        $\Gamma \Entails !B \land !C$। हमें $\Gamma \Entails !B$ दिखाना
        है। अतः मान लें $\mSat{M}{\Gamma}[w]$। चूँकि
        $\Gamma \Entails !B \land !C$, इसलिए
        $\mSat{M}{!B \land !C}[w]$। तब $\mSat{M}{!B}[w]$ भी। इसी प्रकार यदि
        $\Elim\land$ का निष्कर्ष~$!C$ है, तो $\Gamma \Entails !C$।}

  \item !!{derivation} $\Intro{\lor}$ पर समाप्त होती है:
    \iftag{probOr}{अभ्यास।}{मान लें कि पूर्वाधार $!B$ है और~$!B$ पर समाप्त
      होने वाली !!{derivation} की !!{undischarged} मान्यताएँ $\Gamma$ हैं। तब
      $\Gamma \Proves !B$, और आगमन-परिकल्पना से
      $\Gamma \Entails !B$। हमें $\Gamma \Entails !B \lor !C$ दिखाना है।
      मान लें $\mSat{M}{\Gamma}[w]$। चूँकि $\Gamma \Entails !B$,
      $\mSat{M}{!B}[w]$। तब $\mSat{M}{!B \lor !C}[w]$ भी। इसी प्रकार यदि
      पूर्वाधार~$!C$ है, तो $\Gamma \Entails !C$।}
    
  \item !!{derivation}~$\Elim{\lor}$ पर समाप्त होती है:
    \iftag{probOr}{अभ्यास।}{पूर्वाधारों पर समाप्त होने वाली
      !!{derivation}याँ: !!{undischarged} मान्यताओं~$\Gamma$ से
      $!B \lor !C$ की, !!{undischarged} मान्यताओं
      $\Delta_1 \cup \{!B\}$ से $!D$ की, और !!{undischarged} मान्यताओं
      $\Delta_2 \cup \{!C\}$ से $!D$ की हैं। अतः
      $\Gamma \Proves !B \lor !C$, $\Delta_1 \cup \{!B\} \Proves !D$
      और $\Delta_2 \cup \{!C\} \Proves !D$। आगमन-परिकल्पना से
      $\Gamma \Entails !B \lor !C$, $\Delta_1 \cup \{!B\} \Entails !D$
      और $\Delta_2 \cup \{!C\} \Entails !D$। हमें
      $\Gamma \cup \Delta_1 \cup \Delta_2 \Entails !D$ सिद्ध करना है।

    मान लें $\mSat{M}{\Gamma \cup \Delta_1 \cup \Delta_2}[w]$। तब
    $\mSat{M}{\Gamma}[w]$; और चूँकि $\Gamma \Entails !B \lor !C$,
    $\mSat{M}{!B \lor !C}[w]$। $\mSat{M}{}$ की परिभाषा से या तो
    $\mSat{M}{!B}[w]$ या $\mSat{M}{!C}[w]$। अतः मामलों में भेद करें:
    (a)~$\mSat{M}{!B}$[w]। तब $\mSat{M}{\Delta_1 \cup \{!B\}}[w]$।
    चूँकि $\Delta_1 \cup !B \Entails !D$, इसलिए $\mSat{M}{!D}[w]$।
    (b)~$\mSat{M}{!C}[w]$। तब $\mSat{M}{\Delta_2 \cup \{!C\}}[w]$।
    चूँकि $\Delta_2 \cup !C \Entails !D$, इसलिए $\mSat{M}{!D}[w]$। अतः
    दोनों मामलों में $\mSat{M}{!D}[w]$, जैसा हमें दिखाना था।}

  \item !!{derivation} $\Intro{\lif}$ पर समाप्त होकर~$!B\lif !C$ निष्कर्ष
    देती है। तब पूर्वाधार~$!C$ है और उस पूर्वाधार पर समाप्त होने वाली
    !!{derivation} की !!{undischarged} मान्यताएँ
    ~$\Gamma \cup \{!B\}$ हैं। अतः
    $\Gamma \cup \{!B\} \Proves !C$, और आगमन-परिकल्पना से
    $\Gamma \cup \{!B\} \Entails !C$। हमें
    $\Gamma \Entails !B \lif !C$ दिखाना है।

    मान लें $\mSat{M}{\Gamma}[w]$। हमें दिखाना है कि $Rww'$ वाले सभी
    $w'$ के लिए, यदि $\mSat{M}{!B}[w']$ तो $\mSat{M}{!C}[w']$। अतः मान लें
    $Rww'$ और $\mSat{M}{!B}[w']$।
    \olref[sem][rel]{prop:true-monotonic} से
    $\mSat{M}{\Gamma}[w']$। चूँकि $\Gamma \cup \{!B\} \Entails !C$,
    $\mSat{M}{!C}[w']$, जैसा हमें दिखाना था।

  \item !!{derivation} $\Elim{\lif}$ पर समाप्त होती है और निष्कर्ष~$!C$
    है। पूर्वाधार $!B \lif !C$ और $!B$ हैं, जिनकी !!{derivation}याँ
    क्रमशः !!{undischarged} मान्यताओं $\Gamma$, $\Delta$ से हैं। अतः
    $\Gamma \Proves !B \lif !C$ और $\Delta \Proves !B$।
    आगमन-परिकल्पना से $\Gamma \Entails !B \lif !C$ और
    $\Delta \Entails !B$। हमें $\Gamma \cup \Delta \Entails !C$ दिखाना
    है।

    मान लें $\mSat{M}{\Gamma \cup \Delta}[w]$। चूँकि
    $\mSat{M}{\Gamma}[w]$ और $\Gamma \Entails !B \lif !C$,
    $\mSat{M}{!B \lif !C}[w]$। परिभाषा से इसका अर्थ है कि $Rww'$ वाले सभी
    $w'$ के लिए, यदि $\mSat{M}{!B}[w']$ तो $\mSat{M}{!C}[w']$। चूँकि $R$
    स्वपरावर्ती है, $w$ भी उन $w'$ में है जिनके लिए $Rww'$; अर्थात यदि
    $\mSat{M}{!B}[w]$ तो $\mSat{M}{!C}[w]$। चूँकि
    $\mSat{M}{\Delta}[w]$ और $\Delta \Entails !B$,
    $\mSat{M}{!B}[w]$। अतः $\mSat{M}{!C}[w]$, जैसा हमें दिखाना था।

  \item !!{derivation} $\FalseInt$ पर समाप्त होकर~$!A$ निष्कर्ष देती है।
    पूर्वाधार $\lfalse$ है और पूर्वाधार की !!{derivation} की
    !!{undischarged} मान्यताएँ~$\Gamma$ हैं। तब
    $\Gamma \Proves \lfalse$। आगमन-परिकल्पना से
    $\Gamma \Entails \lfalse$। हमें $\Gamma \Entails !A$ दिखाना है।

    हम अप्रत्यक्ष विधि अपनाते हैं। यदि $\Gamma \Entails/ !A$, तो कोई
    प्रतिरूप~$\mModel{M}$ और जगत~$w$ ऐसा है कि
    $\mSat{M}{\Gamma}[w]$ और $\mSat/{M}{!A}[w]$। चूँकि
    $\Gamma \Entails \lfalse$, इसलिए $\mSat{M}{\lfalse}[w]$। किंतु यह
    असंभव है, क्योंकि परिभाषा से $\mSat/{M}{\lfalse}[w]$। अतः
    $\Gamma \Entails !A$।
  \item !!{derivation} $\Intro\lnot$ पर समाप्त होती है: अभ्यास।
  \item !!{derivation} $\Elim\lnot$ पर समाप्त होती है: अभ्यास।
  \end{enumerate}
\end{proof}

\begin{prob}
  \olref[int][sc][snd]{thm:soundness} का प्रमाण पूरा कीजिए।
  $\Intro\lnot$ और $\Elim\lnot$ के मामलों के लिए
  \olref[int][sem][rel]{defn:true-at-w} में $\mSat{M}{\lnot !A}[w]$ की
  परिभाषा का उपयोग कीजिए; अर्थात $\lnot !A$ को $!A \lif \lfalse$ द्वारा
  परिभाषित न मानिए।
\end{prob}

\begin{prob}
  दिखाइए कि निम्न !!{formula} अंतःप्रज्ञावादी तर्कशास्त्र में !!{derivable}
  नहीं हैं:
  \begin{enumerate}
    \item $(!A \lif !B) \lor (!B \lif !A)$
    \item $(\lnot\lnot !A \lif !A) \lif (!A \lor \lnot !A)$
    \item $(!A \lif !B \lor !C) \lif \bigl((!A \lif !B)\lor(!A \lif !C)\bigr)$
  \end{enumerate}
\end{prob}

\end{document}
